\documentclass[11pt,letterpaper]{article}
\usepackage[T1]{fontenc}
\usepackage{lmodern,microtype,amsmath,booktabs,tabularx,array,graphicx,longtable}
\usepackage[margin=2.5cm]{geometry}
\usepackage{xcolor,listings,xurl}
\usepackage[colorlinks=true,urlcolor=blue,linkcolor=blue]{hyperref}
\lstset{language=C++,basicstyle=\small\ttfamily,columns=fullflexible,
  keepspaces=true,breaklines=true,showstringspaces=false,frame=single}
\newcommand{\ffmt}{\texttt{ffmt}}
\newcommand{\fmtlib}{\texttt{\{fmt\}}}
\setlength{\parskip}{0.3em}
\setlength{\parindent}{0pt}
\emergencystretch=2em
% Generated by tools/paper_tables.py; edit the evidence or generator.
\newcommand{\qualificationtable}{%
{\small
\begin{longtable}{llllrr}
\toprule
System & Route & Opt. & Source & Exact/stdout & Exact/numeric \\
\midrule
\endhead
Rome & host-buffer & O0 & 21f9f5ef & 54/54 & 8216/8216 \\
Rome & host-buffer & O2 & 21f9f5ef & 54/54 & 8216/8216 \\
Rome & host-specifiers & O0 & 21f9f5ef & 54/54 & 8216/8216 \\
Rome & host-specifiers & O2 & 21f9f5ef & 54/54 & 8216/8216 \\
Rome & kokkos-serial & O0 & 21f9f5ef & 54/54 & 8216/8216 \\
Rome & kokkos-serial & O2 & 21f9f5ef & 54/54 & 8216/8216 \\
GH200 & host-buffer & O0 & 21f9f5ef & 54/54 & 8216/8216 \\
GH200 & host-buffer & O2 & 21f9f5ef & 54/54 & 8216/8216 \\
GH200 & host-specifiers & O0 & 21f9f5ef & 54/54 & 8216/8216 \\
GH200 & host-specifiers & O2 & 21f9f5ef & 54/54 & 8216/8216 \\
GH200 & cuda & O0 & 21f9f5ef & 54/54 & 8216/8216 \\
GH200 & cuda & O2 & 21f9f5ef & 54/54 & 8216/8216 \\
GH200 & nvcc & O0 & 21f9f5ef & 54/54 & 8216/8216 \\
GH200 & nvcc & O2 & 21f9f5ef & 54/54 & 8216/8216 \\
MI300X & hip & O0 & 21f9f5ef & 54/54 & 8216/8216 \\
MI300X & hip & O2 & 21f9f5ef & 54/54 & 8216/8216 \\
MI300X & kokkos-hip & O0 & 21f9f5ef & 54/54 & 8216/8216 \\
MI300X & kokkos-hip & O2 & 21f9f5ef & 54/54 & 8216/8216 \\
MI300X & acpp & O0 & 21f9f5ef & 52/54 & 8216/8216 \\
A40 & cuda & O0 & 21f9f5ef & 54/54 & 8216/8216 \\
A40 & cuda & O2 & 21f9f5ef & 54/54 & 8216/8216 \\
A40 & nvcc & O0 & 21f9f5ef & 54/54 & 8216/8216 \\
A40 & nvcc & O2 & 21f9f5ef & 54/54 & 8216/8216 \\
Intel Max & sycl & O0 & 21f9f5ef & 53/54 & 8216/8216 \\
Intel Max & sycl & O2 & 21f9f5ef & 53/54 & 8216/8216 \\
Intel Max & kokkos-sycl & O0 & 21f9f5ef & 53/54 & 8216/8216 \\
Intel Max & kokkos-sycl & O2 & 21f9f5ef & 53/54 & 8216/8216 \\
Intel Max & omp-spirv & O0 & 21f9f5ef & 53/54 & 8216/8216 \\
A40 & kokkos-cuda & O0 & 21f9f5ef & 8/8 & 8216/8216 \\
A40 & kokkos-cuda & O2 & 21f9f5ef & 8/8 & 8216/8216 \\
A40 & omp-cuda & O0 & 21f9f5ef & 8/8 & execution\_failure \\
A40 & omp-cuda & O2 & 21f9f5ef & 8/8 & 8216/8216 \\
A40 & acpp & O0 & 21f9f5ef & 7/8 & 8216/8216 \\
A40 & acpp & O2 & 21f9f5ef & 7/8 & 8216/8216 \\
MI300X & omp-hip & O0 & 98dba67c & 8/8 & execution\_failure \\
MI300X & omp-hip & O2 & 98dba67c & 6/8 & 8216/8216 \\
MI300X & acpp & O0 & 98dba67c & 7/8 & 8216/8216 \\
MI300X & acpp & O2 & 98dba67c & 7/8 & 8216/8216 \\
Intel Max & omp-spirv & O2 & 98dba67c & 4/8 & 8216/8216 \\
Intel Max & chipstar & O0 & 98dba67c & 8/8 & 8216/8216 \\
Intel Max & chipstar & O2 & 98dba67c & 8/8 & 8216/8216 \\
\bottomrule
\end{longtable}
}}
\newcommand{\benchmarktable}{%
{\scriptsize
\begin{longtable}{llllllr}
\toprule
System & Route & Engine & Data & Sink & Headers & ns/call \\
\midrule
\endhead
Intel Max & chipstar & ffmt & float \{\} & memory & 334ca65c & 265705.39 \\
Intel Max & sycl & ffmt & float .8e & memory & 332a7a7f & 30529.35 \\
Intel Max & sycl & ffmt & float \{\} & memory & 334ca65c & 261119.40 \\
Intel Max & sycl & ffmt & float .8e & print & 332a7a7f & 69285.14 \\
Intel Max & sycl & ffmt & integer & memory & 332a7a7f & 11340.19 \\
Intel Max & sycl & ffmt & integer & print & 332a7a7f & 44311.77 \\
Intel Max & sycl & printf & float .8e & print & 332a7a7f & 7080.85 \\
Intel Max & sycl & printf & integer & print & 332a7a7f & 6560.48 \\
Arm/GH200 & host-buffer & ffmt & float .8e & memory & 332a7a7f & 122.77 \\
Arm/GH200 & host-buffer & ffmt & float .8e & print & 332a7a7f & 130.50 \\
Arm/GH200 & host-buffer & ffmt & integer & memory & 332a7a7f & 98.54 \\
Arm/GH200 & host-buffer & ffmt & integer & print & 332a7a7f & 97.88 \\
Arm/GH200 & host-buffer & fmt & float .8e & memory & 332a7a7f & 53.84 \\
Arm/GH200 & host-buffer & fmt & integer & memory & 332a7a7f & 48.94 \\
Arm/GH200 & host-buffer & printf & float .8e & print & 332a7a7f & 135.75 \\
Arm/GH200 & host-buffer & printf & integer & print & 332a7a7f & 54.00 \\
Arm/GH200 & host-specifiers & ffmt & float .8e & memory & 332a7a7f & 133.48 \\
Arm/GH200 & host-specifiers & ffmt & float .8e & print & 332a7a7f & 131.50 \\
Arm/GH200 & host-specifiers & ffmt & integer & memory & 332a7a7f & 95.59 \\
Arm/GH200 & host-specifiers & ffmt & integer & print & 332a7a7f & 101.50 \\
Arm/GH200 & host-specifiers & fmt & float .8e & memory & 332a7a7f & 53.00 \\
Arm/GH200 & host-specifiers & fmt & integer & memory & 332a7a7f & 42.45 \\
Arm/GH200 & host-specifiers & printf & float .8e & print & 332a7a7f & 134.50 \\
Arm/GH200 & host-specifiers & printf & integer & print & 332a7a7f & 55.50 \\
A40 & nvcc & ffmt & float .8e & memory & 332a7a7f & 6726.67 \\
A40 & nvcc & ffmt & float .8e & print & 332a7a7f & 222867.13 \\
A40 & nvcc & ffmt & integer & memory & 332a7a7f & 4096.61 \\
A40 & nvcc & ffmt & integer & print & 332a7a7f & 138106.19 \\
A40 & nvcc & printf & float .8e & print & 332a7a7f & 29305.85 \\
A40 & nvcc & printf & integer & print & 332a7a7f & 27100.23 \\
MI300X & hip & ffmt & float .8e & memory & 332a7a7f & 17375.98 \\
MI300X & hip & ffmt & float .8e & print & 332a7a7f & 659792.50 \\
MI300X & hip & ffmt & integer & memory & 332a7a7f & 10411.69 \\
MI300X & hip & ffmt & integer & print & 332a7a7f & 499450.78 \\
MI300X & hip & printf & float .8e & print & 332a7a7f & 50579.27 \\
MI300X & hip & printf & integer & print & 332a7a7f & 50848.38 \\
Rome & host-buffer & ffmt & float .8e & memory & 332a7a7f & 166.45 \\
Rome & host-buffer & ffmt & float .8e & print & 332a7a7f & 212.81 \\
Rome & host-buffer & ffmt & integer & memory & 332a7a7f & 107.81 \\
Rome & host-buffer & ffmt & integer & print & 332a7a7f & 143.12 \\
Rome & host-buffer & fmt & float .8e & memory & 332a7a7f & 93.98 \\
Rome & host-buffer & fmt & integer & memory & 332a7a7f & 77.42 \\
Rome & host-buffer & printf & float .8e & print & 332a7a7f & 315.94 \\
Rome & host-buffer & printf & integer & print & 332a7a7f & 110.62 \\
Rome & host-specifiers & ffmt & float .8e & memory & e73f2c71 & 188.89 \\
Rome & host-specifiers & ffmt & float .8e & print & e73f2c71 & 199.61 \\
Rome & host-specifiers & ffmt & integer & memory & e73f2c71 & 110.25 \\
Rome & host-specifiers & ffmt & integer & print & e73f2c71 & 130.86 \\
Rome & host-specifiers & fmt & float .8e & memory & e73f2c71 & 90.43 \\
Rome & host-specifiers & fmt & integer & memory & e73f2c71 & 73.66 \\
Rome & host-specifiers & printf & float .8e & print & e73f2c71 & 235.16 \\
Rome & host-specifiers & printf & integer & print & e73f2c71 & 102.34 \\
\bottomrule
\end{longtable}
}}
\newcommand{\coveragetable}{%
\begin{tabular}{lrrr}
\toprule
Metric & Covered & Instrumented & Percent \\
\midrule
Regions & 3244 & 3519 & 92.19 \\
Executable lines & 3445 & 3740 & 92.11 \\
Branch edges & 2119 & 2540 & 83.43 \\
\bottomrule
\end{tabular}
}
\newcommand{\resourcetable}{%
{\small
\begin{longtable}{llrlrrr}
\toprule
GPU target & Route & Stage & Sink & Calls & Regs/VGPR & Frame/scratch B \\
\midrule
\endhead
gfx942 & hip & 128 & memory & 1 & 97 & 1440 \\
gfx942 & hip & 128 & memory & 2 & 64 & 1520 \\
gfx942 & hip & 128 & print & 1 & 124 & 1600 \\
gfx942 & hip & 128 & print & 2 & 64 & 1744 \\
gfx942 & hip & 512 & memory & 1 & 97 & 1440 \\
gfx942 & hip & 512 & memory & 2 & 64 & 1520 \\
gfx942 & hip & 512 & print & 1 & 124 & 1984 \\
gfx942 & hip & 512 & print & 2 & 64 & 2128 \\
sm\_86 & nvcc & 128 & memory & 1 & 120 & 1392 \\
sm\_86 & nvcc & 128 & memory & 2 & 120 & 1400 \\
sm\_86 & nvcc & 128 & print & 1 & 70 & 1568 \\
sm\_86 & nvcc & 128 & print & 2 & 103 & 1664 \\
sm\_86 & nvcc & 512 & memory & 1 & 120 & 1392 \\
sm\_86 & nvcc & 512 & memory & 2 & 120 & 1400 \\
sm\_86 & nvcc & 512 & print & 1 & 70 & 1952 \\
sm\_86 & nvcc & 512 & print & 2 & 104 & 2432 \\
\bottomrule
\end{longtable}
}}
\newcommand{\measuredgroups}{768}

\title{\texttt{ffmt}: Checked C++ Formatting in Accelerator Kernels}
\author{Working fork development notes}
\date{22 October 2026 --- working revision 3}
\begin{document}
\maketitle
\begin{abstract}
Readable diagnostics inside accelerator kernels require both device-callable
formatting and an output transport with known behavior. We describe a C++20
implementation that shares one checked formatting engine across host, CUDA,
HIP, SYCL and OpenMP target code. Output iterators provide bounded memory
formatting, while direct kernel printing streams through fixed-size staging
storage without a required host logging service. An independent, pinned
\fmtlib\ build supplies exact expected bytes. Qualification separates stdout
agreement, device-memory numeric agreement, invalid-format diagnostics, runtime
controls, source coverage and performance. We report representative JLSE runs
on x86\_64, Arm and NVIDIA, AMD and Intel accelerators, including Kokkos routes.
We retain failures and limitations rather than counting compilation or skipped
execution as conformance. The experiments establish agreement for the sampled
corpus and captured software stacks; they do not establish exhaustive fmt or
strict freestanding conformance, or performance superiority.
\end{abstract}

\section{Scope and evidence}

\ffmt\ lets a programmer write brace-style formatting inside a kernel. A
successful implementation must validate its syntax and argument types,
construct the intended representation, and deliver the resulting bytes.
These obligations are tested separately. The implemented library has no
mandatory host initialization or drain API; normal runtime synchronization
still applies to the kernel and its output.

The original implementation is pinned at
\texttt{3fe258a87675c997667c899c2b0f0cad405aad5a}. The oracle is fmt commit
\texttt{0fc31633e7d53ac895ae0a688b83b7d44f0b0c1f}; both repositories and
the working fork are identified in the artifact\cite{ffmt-source,fmt-source}.
The supplied 21 October outline is preserved as \texttt{outline-v2.tex}.
Its two-path architecture and unexecuted-test ledger describe that historical
review. This revision describes implemented changes and captured experiments.
Tables are generated from selected JSON evidence rather than filled with
estimated pass counts.

The supported surface targets device-suitable fmt behavior. It includes
runtime and wide formats, positional and named arguments, dynamic
specifications, custom formatters, ranges/tuples/join, locale-independent
chrono, selected standard type adapters, explicit ANSI styling and typed
percent formatting. Filesystem, stream, locale and timezone-database services,
RTTI, thread identifiers and exception traversal are outside this device
surface. A standard type adapter does not make every implementation of that
type device callable.

\section{Interface and shared engine}

\begin{lstlisting}
#include <ffmt/base.hpp>
// Within a supported host function or kernel:
FFMT_PRINTLN("id={0:04d} value={1:.6e} again={0}", id, value);

char buffer[32];
auto result = ffmt::format_to_n(buffer, sizeof buffer,
                               ffmt::runtime(pattern), value);
if (!result) { /* inspect result.error */ }
\end{lstlisting}

Literal-format constructors validate the syntax and the selected formatter
at compilation. Runtime text uses the same parser and returns an error.
\texttt{format\_to\_n} reports the number of code units stored and, on success,
the complete required size. It never writes beyond the supplied count and
does not append a terminator. A short destination is successful bounded
formatting. Runtime errors may leave a prefix; the required count is not
specified on failure. \texttt{try\_print} and \texttt{try\_println} return
checked stdout status. The convenience print calls discard it.

The implementation replaces format-to-printf translation with a shared
type-directed parser and sink. The same format traversal can count output,
write an output iterator or feed the stdout staging sink. Runtime widths and
precisions are checked for their type and value. Format text and field counts
have no separate fixed descriptor-array limit. The default 128-byte staging
size bounds an emission chunk, not the complete output length. Independent
calls may interleave between chunks; atomic whole-message printing is not
part of the contract.

\begin{center}
\fbox{\parbox{0.9\linewidth}{\centering
literal validation / runtime checks $\longrightarrow$ shared format traversal
$\longrightarrow$ typed formatters\\[0.4em]
counting or bounded iterator sink $\quad|\quad$ staging sink
$\longrightarrow$ backend output}}
\end{center}

Formatter contexts expose parsing and output operations without imposing
allocation or exception unwinding in the core. Existing static custom
formatters returning nested \texttt{formatted} values remain accepted.
New custom types can derive from another formatter and reuse its parser.
Named fields refer to their original values; argument expressions are not
repeated to implement repeated positional fields. The API does not create
new left-to-right argument-evaluation guarantees.

Padding a family must not accidentally invoke a mutating formatter twice.
Bounded pointer destinations render once and move the retained prefix in place;
left-aligned streaming also renders once. Right/center alignment on an
irreversible sink requires a body proven replayable by the library. Unknown
custom formatters and consuming ranges return \texttt{unsupported\_type} before
rendering that body. The implementation does not introduce an arbitrary
temporary-output cap to hide this restriction.

\section{Conversion and transport}

Integer formatting preserves signed magnitude for nondecimal brace formats
and checks the type-specific specification. Decimal floating conversion uses
integer representations of IEEE values. The ordinary binary32/binary64
shortest path incorporates Dragonbox material adapted through
fmt\cite{dragonbox,fmt-source}. An exact decimal interval implementation supplies
a portable shortest path for SPIR-V. Extended host representations have
dedicated conversion rather than narrowing to double. Native extended-format
tests distinguish x87 precision from Arm binary128. A separate native
\texttt{\_\_float128} test uses a transparent oracle formatter adapter because
the pinned fmt public type mapper rejects that raw extension type.

Fixed and scientific conversion, hexadecimal conversion, signed zero,
subnormals and special encodings are exercised separately. Large requested
precision is emitted through the sink, so increasing precision does not
require an output-sized temporary. This does not imply that stack use is
constant across types or compiler instantiations: numeric scratch state,
inlining, spills and device runtime call frames remain relevant.

Narrow memory output preserves bytes, including embedded NUL. Wide memory
formatting matches the oracle's code units. Wide stdout is encoded as UTF-8;
invalid scalar values and unmatched UTF-16 surrogate halves produce an error.
The Unicode tests cover selected width, precision, fill and debug-escaping
behavior. They are not a complete Unicode conformance suite.

The host transport uses counted byte output. Native accelerator transports
use constant percent-conversion strings to deliver already formatted bytes;
formatted payload is not reinterpreted as a percent format. This avoids a
dependency on runtime-produced string pointers being accepted by a device
printf implementation. Backend output remains a runtime dependency.

The measured Intel SPIR-V runtime drops NUL even through a literal character
conversion. The default implementation detects NUL-containing chunks and
returns \texttt{transport\_error}; it does not report successful complete
printing. Memory formatting is independent of that limitation. An opt-in
macro permits a separately validated binary-safe transport. Runtime output
buffer exhaustion, concurrent interleaving and kernel synchronization are
distinct from formatting correctness.

AdaptiveCpp's supported SSCP print builtin accepts null-terminated text. The
measured route therefore also returns \texttt{transport\_error} for a chunk
containing NUL. An independent builtin probe emits only the prefix before NUL;
memory output preserves all bytes. Direct variadic printf failed even for an
ASCII control in that stack, so it supplies no alternative binary-safe path.
This limitation belongs to the measured SSCP transport, not native HIP printf.

\section{Independent conformance method}

The oracle source and its implementation object are compiled for the host.
Generated target programs include ffmt and backend headers, never fmt.
Expected output is captured without trimming it to an ffmt capacity. Each
stdout case is invoked in a separate process, and raw byte comparison retains
empty output, newlines, invalid UTF-8 and NUL. Payload cannot imitate a harness
control record. Harness regressions deliberately test truncation, extra and
duplicate bytes, reordering, missing records and unrelated compiler errors.

Every route first builds and runs a positive control and reports the actual
selected execution space or device. CUDA/HIP launches check runtime errors;
SYCL selects a GPU and waits with exception propagation; OpenMP checks
\texttt{omp\_is\_initial\_device} with mandatory offload. Kokkos verifies the
expected default execution space in stdout and independent memory controls.
An unavailable compiler, missing device or rejected control is not a pass.

Compilation failures are recursively isolated from shared source shards.
Execution-failing shards are also reduced before assigning per-case outcomes,
because a shared switch can change stack allocation or compiler behavior.
The failed combined binaries and diagnostics remain in the manifest. This
distinguishes individual formatting failures from code-generation or resource
failures affecting a larger translation unit.

The numeric memory corpus loads 4,096 seeded binary32 patterns and 4,096
binary64 patterns, plus 24 explicit boundary encodings, from device memory.
Each of the 8,216 input records combines default, fixed, scientific, general
and hexadecimal formatting. Results, status and required/written counts are
copied back and serialized with checked framing. Thus these counts denote
input records, each containing five format checks; they are not 8,216 stdout
messages. This suite runs independently even when a stdout control fails.

Invalid-literal testing compares intended rejection in both libraries.
Positive and per-file syntax controls prevent missing headers or malformed
fixtures from satisfying the negative test. Unrelated diagnostics and compiler
crashes have separate statuses. Runtime-invalid text is tested through checked
status APIs. Optional compiler types use an explicit unsupported status.

\section{Representative JLSE qualification}

The selected systems are Rome x86\_64, GH200 Arm/CUDA, NVIDIA A40 CUDA,
MI300X HIP and Intel Max SYCL/SPIR-V. The H100 queue profile used an explicitly
recorded A40 alternative because the primary queue was occupied. The final
qualification does not establish H100 execution; the original baseline separately
records that device. These are Cobalt allocations,
not login-node GPU executions. Profiles retain compiler flags, device identity,
driver/runtime discovery, source hashes and commands. Optional toolchains and
runtimes are installed under ignored per-system software prefixes with pinned
recipes. Kokkos, AdaptiveCpp and chipStar remain distinct routes rather than
aliases for native CUDA/HIP/SYCL support.

\qualificationtable
Selected captured qualification configurations. Stdout and numeric
denominators are separate. A smaller exact count and every non-pass numeric
status retain a gap; a table row is not a blanket backend-support claim.
Per-case evidence, source snapshots and diagnostics are in the companion JSON.

Strict configurations use \texttt{-fno-fast-math} at both optimization levels.
Host sanitizers and strict warnings exercise bounded canaries, parser errors,
wide data and custom protocols. Fast-math is a separate assessment: a
formatter cannot reconstruct information destroyed by earlier user arithmetic.
JIT startup limits are recorded. One Intel optimized numeric binary initially
timed out at 120 seconds; the identical binary subsequently matched all 8,216
records with a longer startup allowance. That timeout is not evidence of a
numeric mismatch.

\subsection{Source coverage and coverage gaps}

\begin{center}
\coveragetable
\end{center}

Source coverage unions executable positions across independently profiled
binaries and template instantiations at snapshot \texttt{0c2cb76}, including
the final counted-view validation repair. It reports regions, executable lines
and true/false branch edges for all instrumented ffmt headers. The denominator
excludes uninstantiated templates, compile-time evaluation and device-only
code. It therefore cannot establish hardware coverage or completeness of the
fmt feature surface. The functional corpus also samples rather than exhausts
format strings, types, floating encodings and Unicode text.

\section{Measurements and resources}

Formatting into preallocated memory and direct printing are measured
separately. Inputs and outputs remain observable through exact verification
and checksums. The formatting benchmark does not charge only one library for
allocating a result string. Host fmt is compared on host memory formatting;
it is not represented as a GPU implementation. Native printf is compared
only for semantically matched integer/fixed-scientific output specifications.
The sample does not compare default shortest conversion with unqualified
percent-g formatting.

Each measured group has five warm-up batches and thirty measured batches.
The host wall timer for accelerators includes launch and synchronization,
with data copies and checksums outside the timed interval. Normalizing batch
time by calls gives an amortized figure, not individual GPU-call latency.
Intel SYCL additionally records device-event batch times and explicit
executable-kernel-bundle preparation with persistent SYCL cache disabled.
Backend driver caches can still affect preparation; these measurements are
not a claim of hardware-cache-cold compilation. Historical first-batch timing
after an earlier verification process is labeled first batch, not cold JIT.

\benchmarktable
Representative measured groups: one call, one translation unit and
128-byte staging. Values are median batch time divided by calls. Host and
device execution models differ; cross-system rows are not speedup claims.
The full artifact contains \measuredgroups\ measured groups, including other
capacities, two calls and multiple translation units. For overlapping groups,
the table selects the latest captured measured run. The header fingerprint
identifies its exact source map in \texttt{paper-selection.json}; historical
measurements are not relabeled as the final revision.

Resource probes separately instantiate integer/floating formats, memory/stdout
sinks and one/two calls at staging sizes 16, 64, 128 and 512. Fresh output
directories prevent object reuse for compilation timing, although filesystem
caches are not evicted. Objects, final executables, tool reports and hashes
are retained. Compiler static-frame reports exclude cumulative call depth;
device stack frames, spills, register counts, read-only sections and executable
size are different quantities. A source staging array is not evidence that
all of it occupies shared local memory, nor does an object-size change alone
measure runtime occupancy. Unsupported resource-report fields remain absent.

\resourcetable
These compile-only shortest-format probes target the named GPU architecture;
they do not require that GPU on the compilation node. They report CUDA compiler stack-frame bytes
or AMD scratch bytes per lane, respectively; those columns are not equivalent
cross-vendor measures. Registers are CUDA registers or AMD VGPRs. Device-image
symbol inspection also distinguishes read-only formatting tables before and after
linking. The captured CUDA multi-translation-unit build retains two sets of
four tables (4,096 bytes), versus one set (2,048 bytes) in a single translation
unit. Multiple call sites within one translation unit do not independently
duplicate those tables in this probe.

\section{Limits and reproducibility}

Results belong to the captured source and runtime stacks, not to every
release sharing a backend name. Compile-only minimized probes are labeled
accordingly. The artifact records failures before and after fixes, including
device-callability of tuple construction, nonrecursive chrono lowering,
matched OpenMP compiler/runtime requirements and numeric launch resource
failures. No upstream compiler issue is claimed filed by this experiment.

The stdout conformance probes use one active work item per case. They do not
establish concurrent-message behavior or printf FIFO saturation limits. The
printing measurements use bounded experiment volumes. No separate no-op launch
baseline, \texttt{std::format} or \texttt{FMT\_COMPILE} comparison was collected;
the reported host-wall batch timing includes launch overhead rather than
subtracting an estimated baseline.

The matched ROCm OpenMP route has an optimization-dependent failure. A
97-line independent callback reproducer with no fmt or ffmt dependency aborts
at a null GPU address at \texttt{-O0}, while the identical source passes 8,216
records at \texttt{-O2}. A host OpenMP control passes under address and
undefined-behavior sanitizers. This isolates a compiler/runtime boundary for
that captured stack; it does not establish that all OpenMP failures have the
same cause. Source, exact commands and outcomes accompany the artifact.

The NVIDIA OpenMP numeric probe also reports an illegal access at
\texttt{-O0}, while its optimized numeric run passes. Focused OpenMP family
builds can exceed the recorded compiler timeout; an optimized Intel padding
binary additionally fails module creation with an unresolved overflow
intrinsic. These outcomes remain gaps rather than exact-output passes. The
independent AMD reproducer does not by itself explain the other failures.

The library is header-only and exports a CMake interface target. Make and
CTest run the focused and independent suites; the JLSE helper submits
allocation-owned jobs. The README and evidence guide provide exact commands.
Raw installations, logs and binaries are untracked; compact evidence and
pinned recipes are tracked. The fork's stacked PRs separate harness,
engine, family, portability, qualification and measurement changes.
Further systems can use the same oracle and record schema without converting
a skipped or unexecuted configuration into a passing result.

\begin{thebibliography}{9}
\bibitem{ffmt-source} T. Applencourt, \emph{ffmt}, original source and working
fork. \url{https://github.com/TApplencourt/ffmt},
\url{https://github.com/nscottnichols/ffmt}. Immutable revisions and per-file
hashes are recorded in the evidence manifests.
\bibitem{fmt-source} V. Zverovich and contributors, \emph{fmt}, oracle source.
\url{https://github.com/fmtlib/fmt/tree/0fc31633e7d53ac895ae0a688b83b7d44f0b0c1f}.
\bibitem{dragonbox} J. Jeon and contributors, \emph{Dragonbox}.
\url{https://github.com/jk-jeon/dragonbox}. Adaptation and component licenses
are identified in the artifact's NOTICE and LICENSES directory.
\end{thebibliography}
\end{document}
